\chapter{Выбор инструментов и технологий для реализации системы мониторинга} \label{ch2}
	
% не рекомендуется использовать отдельную section <<введение>> после лета 2020 года
%\section{Введение} \label{ch2:intro}

\section{Система сбора и хранения метрик} \label{ch2:sec1}

Как показано в предыдущей главе, система сбора и хранения метрик является основной составляющей разрабатываемого решения мониторинга. 

На основании требований, сформулированных в предыдущей главе, определим основные критерии, предъявляемые к инструменту автоматизированного сбора и хранения метрик:

\begin{itemize}
	\item возможность регулярного автоматизированного сбора метрик из различных источников, включая системные представления Greenplum и средства мониторинга операционной системы;
	
	\item поддержка хранения временных рядов с сохранением истории изменений показателей;
	
	\item наличие встроенных механизмов агрегации, фильтрации и обработки поступающих метрик;
	
	\item совместимость с внешними системами визуализации и оповещения;
	
	\item простота настройки и администрирования;
	
	\item отсутствие необходимости использования коммерческой лицензии.
\end{itemize}

Среди наиболее распространенных свободно распространяемых инструментов, удовлетворяющих указанным требованиям, можно выделить Zabbix и Prometheus. Оба решения широко применяются для построения систем мониторинга и позволяют реализовать сбор, хранение и анализ метрик, а также интегрируются со средствами визуализации и оповещения.

Рассмотрим основные компоненты архитектуры Zabbix.

Zabbix Server -- центральный элемент системы, который выполняет прием данных от подключенных узлов, их обработку, сохранение в базе данных. Кроме того, сервер отвечает за формирование уведомлений и выполнение настроенных сценариев оповещения.

Zabbix Agent обеспечивает сбор информации о состоянии отслеживаемых хостов и передачу собранных значений на сервер мониторинга.

Zabbix Frontend предоставляет веб-интерфейс управления мониторингом. Через него выполняется настройка, и конфигурирование системы, а также построение дашбордов.

Для хранения накопленных метрик в Zabbix используется внешняя реляционная база данных. В качестве системы хранения могут использоваться PostgreSQL, MySQL, Oracle и другие поддерживаемые СУБД.

Также в системе реализован встроенный механизм оповещений, поддерживающий отправку уведомлений по различным каналам связи.

Таким образом, Zabbix представляет собой универсальную многофункциональную платформу инфраструктурного мониторинга.

Несмотря на широкие функциональные возможности Zabbix, его использование в рамках разрабатываемой системы мониторинга СУБД Greenplum имеет ряд существенных ограничений.

Основным недостатком является отсутствие готовых специализированных плагинов для работы с СУБД Greenplum. Для получения диагностических показателей в Zabbix потребовалась бы разработка алгоритма сбора метрик, специфических для Greenplum. При расширении списка собираемых метрик это приводит к существенному увеличению трудоёмкости сопровождения и усложняет администрирование системы.

Еще одной особенностю инструмента является способ хранения данных. В Zabbix -- используется внешняя реляционная база данных, что ограничивает возможности оперативной аналитической обработки временных рядов. В частности, в Zabbix отсутствует специализированный язык запросов, аналогичный PromQL, предназначенный для гибкой агрегации, фильтрации и вычисления производных метрик.

Перейдем к рассмотрению Prometheus. Prometheus представляет собой комплексное решение для сбора, хранения и анализа метрик, разработанное специально для задач инфраструктурного и сервисного мониторинга.

Архитектурно Prometheus включает несколько основных компонентов.

Центральным элементом является Prometheus Server, который выполняет периодический опрос источников метрик по протоколу HTTP, сохраняет полученные значения во встроенное хранилище временных рядов TSDB (Time Series Database), а также обеспечивает выполнение запросов к накопленным данным.

Источниками метрик в архитектуре Prometheus выступают экспортеры -- вспомогательные программные модули, публикующие показатели работы наблюдаемой системы в формате, понятном Prometheus. Сервер мониторинга регулярно обращается к экспортерам по pull-модели, считывая актуальные значения показателей. Такой подход упрощает управление процессом сбора данных и позволяет легко масштабировать систему при добавлении новых контролируемых узлов.

Основой модели данных в Prometheus является метрика, представляющая собой временной ряд. Каждая метрика определяется как набор значений, изменяющихся во времени, и хранится в виде последовательности пар «временная метка -- значение».

Особенностью модели данных Prometheus является использование меток (labels), представляющих собой набор пар «ключ -- значение», ассоциированных с каждой метрикой. Метки позволяют однозначно идентифицировать временной ряд и задавать дополнительные измерения, такие как имя хоста, имя базы данных, идентификатор сегмента или другие характеристики контролируемого объекта.

Таким образом, каждая метрика в Prometheus определяется уникальной комбинацией имени метрики и набора меток. Такой подход обеспечивает высокую гибкость при организации мониторинга, позволяя агрегировать, фильтровать и анализировать данные по различным срезам.

Нотация метрики в Prometheus иммеет следующий вид:\\
 \texttt{<metric name>\{<label name>="<label value>", ...\}} 
 
Пример метрики \texttt{pg\_locks\_count}:\\
\texttt{pg\_locks\_count\{datname="postgres",mode="accesssharelock"\}} 

Для обработки накопленных временных рядов в Prometheus используется специализированный язык запросов PromQL (Prometheus Query Language). Он предоставляет возможности фильтрации, агрегации, вычисления скоростей прироста и других производных расчетов. Это особенно важно при анализе метрик, когда требуется не только фиксировать абсолютные значения, но и выявлять динамику изменения показателей.




\section{Архитектура разрабатываемой системы мониторинга} \label{ch2:architecture} %название по-русски

В рамках данной работы была разработана система мониторинга производительности СУБД Greenplum, реализованная с использованием контейнеризации на основе Docker. Такой подход позволяет обеспечить изоляцию компонентов, упростить развертывание системы и повысить её воспроизводимость.

Развёртывание системы осуществляется с использованием инструмента Docker Compose, который позволяет описать конфигурацию всех компонентов в одном файле и запускать их в рамках единой инфраструктуры.

В состав разработанной системы входят следующие компоненты:
\begin{itemize}
	\item СУБД Greenplum, используемая в качестве источника данных для мониторинга. В тестовой конфигурации развернут кластер, включающий управляющий узел и четыре сегмента без репликации;
	\item Prometheus — система сбора и хранения метрик, осуществляющая периодический опрос источников данных;
	\item Grafana — инструмент визуализации, используемый для построения дашбордов и анализа метрик;
	\item Alertmanager — компонент, отвечающий за обработку и отправку уведомлений о критических событиях \cite{alertmanager};
	\item Node Exporter — сервис для сбора системных метрик (нагрузка на процессор, использование памяти, дисковая активность) \cite{nodeexporter};
	\item Postgres Exporter — инструмент для получения метрик, связанных с работой СУБД PostgreSQL и Greenplum на уровне базы данных \cite{exporter};
	\item Greenplum Exporter — специализированный компонент, обеспечивающий сбор метрик, специфичных для СУБД Greenplum.
\end{itemize}

Все перечисленные компоненты развёрнуты в отдельных контейнерах и взаимодействуют между собой через единую Docker-сеть. Такой подход обеспечивает модульность системы и упрощает её масштабирование и сопровождение.

Процесс развертывания системы автоматизирован с использованием конфигурационного файла docker-compose.yml и файла переменных окружения .env. Файл .env содержит параметры конфигурации, такие как учетные данные, адреса сервисов и другие настройки, необходимые для корректной работы системы. Перед запуском системы пользователь должен заполнить данный файл, после чего выполнить команду запуска Docker Compose. В результате все компоненты автоматически инициализируются и начинут взаимодействие.

Сбор метрик осуществляется Prometheus, который опрашивает экспортеры по заданным HTTP-эндпоинтам. Полученные данные сохраняются в виде временных рядов и используются для последующего анализа и визуализации.

Визуализация метрик реализована в Grafana. В рамках системы настроено подключение к Prometheus в качестве основного источника данных. Кроме того, реализовано дополнительное прямое подключение Grafana к базе данных Greenplum. Это позволяет выполнять SQL-запросы непосредственно из интерфейса Grafana и отображать результаты в виде табличных виджетов. Такой подход расширяет возможности анализа и позволяет получать более детализированную информацию о состоянии системы.

Alertmanager используется для обработки событий, связанных с превышением пороговых значений метрик. Он получает сигналы от Prometheus и выполняет отправку уведомлений в соответствии с заданными правилами.

Разработанная архитектура обеспечивает централизованный сбор, хранение и визуализацию метрик, а также реализацию механизма оповещения. Использование контейнеризации и автоматизированного развертывания делает систему удобной в использовании и адаптируемой к различным условиям эксплуатации.

%\FloatBarrier % заставить рисунки и другие подвижные (float) элементы остановиться

\section{Организация сбора метрик} \label{ch2:prom}

В разработанной системе мониторинга ключевую роль в процессе сбора и хранения метрик выполняет система Prometheus. Она представляет собой специализированное решение для работы с временными рядами, предназначенное для получения, хранения и обработки метрик в распределённых системах.

Архитектура Prometheus построена по принципу автономности и включает в себя все необходимые компоненты: механизм сбора данных, встроенное хранилище и язык запросов для анализа. Это позволяет использовать систему без привлечения внешних средств хранения и упрощает её развертывание и сопровождение.

Сбор метрик в Prometheus реализован по pull-модели, при которой система самостоятельно инициирует опрос источников данных через заданные интервалы времени. В качестве источников метрик в данной работе используются специализированные сервисы — экспортеры. Они выполняют преобразование внутренних показателей системы в формат, поддерживаемый Prometheus, и предоставляют их через HTTP-интерфейс.

В разработанной системе используются следующие экспортеры:

\begin{itemize}
	\item Node Exporter — обеспечивает сбор системных метрик, включая загрузку процессора, использование оперативной памяти, состояние дисковой подсистемы и сетевую активность;
	\item Postgres Exporter — предоставляет метрики, характеризующие работу базы данных, такие как количество активных подключений, статистика выполнения запросов и общее состояние СУБД;
	\item Greenplum Exporter — дополняет набор метрик за счёт показателей, специфичных для СУБД Greenplum.
\end{itemize}

Каждый экспортер функционирует как отдельный сервис и предоставляет метрики по HTTP-эндпоинту. Prometheus обращается к этим эндпоинтам, извлекает данные и сохраняет их во внутреннем хранилище. Взаимодействие между компонентами осуществляется в рамках одной Docker-сети, что позволяет использовать внутренние имена контейнеров при настройке подключения.

Конфигурация сбора метрик задаётся в файле настроек Prometheus, где определяется список целей (targets), соответствующих адресам экспортеров, а также параметры опроса, включая интервал сбора данных. Это обеспечивает централизованное управление процессом мониторинга и позволяет гибко изменять состав источников метрик.

Полученные данные сохраняются в виде временных рядов, где каждая метрика представлена последовательностью значений, привязанных ко времени. Дополнительно используются метки (labels), позволяющие описывать характеристики метрик и выполнять их агрегацию по различным признакам. Такой подход обеспечивает гибкость при анализе данных и позволяет формировать сложные запросы.

Сбор метрик в рамках данной работы осуществляется с управляющего узла Greenplum, что позволяет получать агрегированную информацию о состоянии системы, активности подключений и использовании ресурсов. Несмотря на ограниченность источника данных, такой подход обеспечивает достаточную информативность для выявления большинства критических ситуаций.

Таким образом, использование Prometheus в сочетании с набором экспортеров позволяет реализовать централизованную систему сбора метрик, обеспечивающую регулярное получение данных, их хранение в виде временных рядов и подготовку к последующему анализу и визуализации.

\section{Визуализация метрик в Grafana} \label{ch2:grafana}

После организации сбора и хранения метрик следующим этапом является их визуализация, позволяющая анализировать состояние системы в наглядной форме. В разработанной системе для этих целей используется Grafana — современный инструмент визуализации данных, предназначенный для работы с временными рядами.

Grafana предоставляет возможность подключения к различным источникам данных и построения интерактивных дашбордов. В рамках данной работы в качестве основного источника данных используется Prometheus, из которого извлекаются собранные метрики. Подключение осуществляется через встроенный механизм Data Source, позволяющий выполнять запросы к системе мониторинга с использованием языка PromQL.

Одним из ключевых преимуществ Grafana является гибкость при построении дашбордов. Пользователь может создавать панели различных типов, включая графики, диаграммы, индикаторы и таблицы. Это позволяет адаптировать отображение данных под конкретные задачи мониторинга и упростить восприятие информации.

В разработанной системе был создан пользовательский дашборд, предназначенный для отображения ключевых показателей состояния СУБД и используемых ресурсов. Дашборд объединяет данные из различных источников и обеспечивает централизованное представление информации о работе системы. При его проектировании основное внимание уделялось наглядности, информативности и удобству восприятия данных.

Структурно дашборд разделён на четыре логических блока, каждый из которых отвечает за отдельный аспект мониторинга системы. Такое разделение позволяет упростить анализ и быстрее ориентироваться в показателях.

На дашборде представлены панели, отражающие текущее состояние системы. В частности, реализованы индикаторы, позволяющие оценить доступность системы и её основных компонентов. Такие элементы обеспечивают быстрое понимание текущего состояния СУБД без необходимости детального анализа графиков.

Отдельный блок дашборда посвящён мониторингу нагрузки на систему. В нём отображаются показатели использования процессора и оперативной памяти, полученные с помощью Node Exporter. Эти метрики позволяют оценить уровень загрузки системы и выявить возможные проблемы, связанные с нехваткой ресурсов.

Также на дашборде представлены метрики, характеризующие работу базы данных. К ним относятся показатели количества активных подключений, активности запросов и общего состояния СУБД. Данные метрики позволяют анализировать поведение пользователей и выявлять аномалии в работе системы.

Дополнительно используются показатели, специфичные для Greenplum, включая информацию о состоянии сегментов и их активности. Это позволяет оценивать общее состояние распределённой системы и контролировать её функционирование.

\begin{figure}[ht!] 
	\center
	\includegraphics [scale=0.13] {/Users/vadimvasilev/Study/8 семестр/НИР/latex/SPbPU-student-thesis-template/my_folder/images//3}
	\caption{Блок "Greenplum Overview"} 
	\label{fig:block1}  
\end{figure}

Первый блок "Greenplum Overview" изображенный на \firef{fig:block1} содержит метрики, отражающие общее состояние системы и её ключевых компонентов:

\begin{itemize}
	\item Cluster status — \lstinline|min(pg_up{})|. Показывает общее состояние кластера. Значение определяется по наименее доступному компоненту.
	
	\item Exporter status — \lstinline|min(up{})|. Отражает доступность компонентов системы мониторинга.
	
	\item Segments up — \lstinline|sum(greenplum_node_segment_status{content!="-1"})|. Показывает количество активных сегментов кластера.
	
	\item Connections count — \lstinline|sum(pg_stat_database_numbackends{job="greenplum"})|. Характеризует текущее количество активных подключений к базе данных.
	
	\item Cluster uptime — \lstinline|greenplum_cluster_uptime|. Отражает время непрерывной работы системы с момента последнего запуска.
	
	\item Cluster size — \lstinline|sum by(job) (pg_database_size_bytes)|. Показывает суммарный объём данных, хранящихся в базе данных.
	
	\item Disk free space — \lstinline|node_filesystem_avail_bytes|. Характеризует объём доступного дискового пространства.
	
	\item CPU usage — \lstinline|(1 - avg by(instance, job)(rate(node_cpu_seconds_total{mode="idle"}[5m]))) * 100|. Отражает уровень загрузки процессора.
	
	\item RAM usage — \lstinline|(1 - (node_memory_MemAvailable_bytes / | \\ \lstinline|node_memory_MemTotal_bytes)) * 100|. Показывает степень использования оперативной памяти.
	
	\item Connections utilization — \lstinline|sum(pg_stat_database_numbackends) / sum(pg_settings_max_connections) * 100|. Характеризует уровень использования доступных подключений.
\end{itemize}

\begin{figure}[ht!] 
	\center
	\includegraphics [scale=0.17] {/Users/vadimvasilev/Study/8 семестр/НИР/latex/SPbPU-student-thesis-template/my_folder/images//4}
	\caption{Блок "Settings"} 
	\label{fig:block2}  
\end{figure}

Второй блок "Settings" изображенный на \firef{fig:block2} содержит параметры конфигурации СУБД Greenplum:

\begin{itemize}
	\item Max WAL size — \lstinline|pg_settings_max_wal_size_bytes{}|. Определяет максимальный размер журнала предзаписи (WAL), при достижении которого инициируется контрольная точка.
	
	\item Work mem — \lstinline|pg_settings_work_mem_bytes{}|. Задаёт объём памяти, выделяемой для выполнения операций сортировки и хеширования в рамках одного запроса.
	
	\item Maintenance mem — \lstinline|pg_settings_maintenance_work_mem_bytes{}|. Определяет объём памяти, используемый для служебных операций, таких как вакуумизация и создание индексов.
	
	\item Shared buffers size — \lstinline|pg_settings_shared_buffers_bytes{}|. Показывает объём общей буферной памяти, используемой СУБД для кэширования данных.
	
	\item Effective cache size — \lstinline|pg_settings_effective_cache_size_bytes{}|. Оценивает объём доступного кэша (включая системный), используемый планировщиком запросов при выборе стратегии выполнения.
	
	\item Random page cost — \lstinline|pg_settings_random_page_cost{}|. Определяет относительную стоимость случайного доступа к страницам данных.
	
	\item Seq page cost — \lstinline|pg_settings_seq_page_cost{}|. Задаёт стоимость последовательного чтения страниц данных.
	
	\item Max parallel workers — \lstinline|pg_settings_max_parallel_workers{}|. Определяет максимальное количество параллельных рабочих процессов, доступных для выполнения запросов.
\end{itemize}

\begin{figure}[ht!] 
	\center
	\includegraphics [scale=0.145] {/Users/vadimvasilev/Study/8 семестр/НИР/latex/SPbPU-student-thesis-template/my_folder/images//5}
	\caption{Блок "Database Metrics"} 
	\label{fig:block3}  
\end{figure}

Третий блок на \firef{fig:block3} включает метрики, характеризующие активность базы данных и эффективность выполнения операций:

\begin{itemize}
	\item TPS (commits) — \lstinline|rate(pg_stat_database_xact_commit{datname=~"$Database"}[$__rate_interval])|. Показывает количество успешно зафиксированных транзакций в секунду.
	
	\item TPS (rollbacks) — \lstinline|rate(pg_stat_database_xact_rollback{datname=~"$Database"}[$__rate_interval])|. Показывает количество откатанных транзакций в секунду.
	
	\item Fetched data — \lstinline|rate(pg_stat_database_tup_fetched{datname=~"$Database"}[$__rate_interval])|. Показывает количество строк, извлечённых из базы данных.
	
	\item Inserted data — \lstinline|rate(pg_stat_database_tup_inserted{datname=~"$Database"}[$__rate_interval])|. Показывает интенсивность операций вставки данных.
	
	\item Deleted data — \lstinline|rate(pg_stat_database_tup_deleted{datname=~"$Database"}[$__rate_interval])|. Показывает количество удалённых строк.
	
	\item Updated data — \lstinline|rate(pg_stat_database_tup_updated{datname=~"$Database"}[$__rate_interval])|. Показывает количество обновлённых строк.
	
	\item Cache hit ratio — \lstinline|pg_stat_database_blks_hit{datname=~"$Database"} / (pg_stat_database_blks_read{datname=~"$Database"} +|  \\ \lstinline|pg_stat_database_blks_hit{datname=~"$Database"}) * 100|. Характеризует долю обращений к данным, обслуженных из кэша.
	
	\item Total cache hit ratio — \lstinline|greenplum_server_database_hit_cache_percent_rate|. Показывает агрегированный показатель эффективности использования кэша.
	
	\item Lock tables — \lstinline|pg_locks_count{datname=~"$Database"} != 0|. Отражает наличие блокировок в базе данных.
	
	\item CPU usage — \lstinline|(1 - avg by(instance, job)(rate(node_cpu_seconds_total{mode="idle"}[$__rate_interval]))) * 100|. Показывает уровень загрузки процессора.
	
	\item Commit success rate — \\ \lstinline|greenplum_server_database_transition_commit_percent_rate|. Показывает долю успешно завершённых транзакций.
	
	\item RAM usage — \lstinline|(1 - (node_memory_MemAvailable_bytes /| \\ \lstinline|node_memory_MemTotal_bytes)) * 100|. Показывает уровень использования оперативной памяти.
\end{itemize}

\begin{figure}[ht!] 
	\center
	\includegraphics [scale=0.15] {/Users/vadimvasilev/Study/8 семестр/НИР/latex/SPbPU-student-thesis-template/my_folder/images//6}
	\caption{Блок "Tables Information"} 
	\label{fig:block4}  
\end{figure}

Четвёртый блок представляет собой табличное представление данных (\firef{fig:block4}), сформированное на основе SQL-запроса к системным каталогам базы данных, как показано в листинге~\ref{lst:tables_info}. Данный блок предназначен для детального анализа состояния таблиц и их активности.

Описание полей табличного представления: 

\begin{itemize}
	\item Schema — схема, к которой относится таблица.
	
	\item Table name — имя таблицы.
	
	\item Relkind — тип объекта (в данном случае отображаются только обычные таблицы).
	
	\item Table size — общий размер таблицы с учётом индексов и вспомогательных данных.
	
	\item Seq scans — количество последовательных сканирований таблицы.
	
	\item Idx scans — количество обращений к таблице с использованием индексов.
	
	\item Inserts — количество операций вставки строк.
	
	\item Updates — количество операций обновления строк.
	
	\item Deletes — количество операций удаления строк.
	
	\item Last vacuum — время последнего ручного выполнения VACUUM.
	
	\item Last autovacuum — время последнего автоматического выполнения VACUUM.
	
	\item Last analyze — время последнего ручного выполнения ANALYZE.
	
	\item Last autoanalyze — время последнего автоматического выполнения ANALYZE.
\end{itemize}

Данный блок позволяет выявлять наиболее нагруженные таблицы, анализировать характер доступа к данным (последовательное или индексное сканирование), а также контролировать актуальность статистики и выполнение операций обслуживания. Это обеспечивает более глубокий уровень мониторинга по сравнению с агрегированными метриками и позволяет принимать обоснованные решения по оптимизации структуры и производительности базы данных.

\section{Организация системы оповещений} \label{ch2:alerts}

Важным элементом системы мониторинга является механизм оповещений, позволяющий своевременно реагировать на отклонения в работе СУБД. В разработанной системе данная функциональность реализована с использованием компонента Alertmanager, входящего в экосистему Prometheus.

Alertmanager предназначен для обработки уведомлений, формируемых системой мониторинга на основе заданных правил. Он обеспечивает агрегацию, фильтрацию и маршрутизацию оповещений, а также предотвращает дублирование сообщений при множественных срабатываниях одного и того же инцидента. Дополнительно реализуется механизм группировки событий, позволяющий объединять связанные алерты в единое уведомление и тем самым снижать информационную нагрузку на пользователя.

Формирование оповещений осуществляется на стороне Prometheus с использованием правил (alerting rules), которые задаются в конфигурационном файле. Каждое правило описывает условие срабатывания на основе выражений языка PromQL. При выполнении заданного условия формируется событие, которое передаётся в Alertmanager для дальнейшей обработки.

В рамках разработанной системы все правила оповещений объединены в единый набор \textit{greenplum-alerts} изображенный на \firef{fig:allalerts}, включающий 11 правил, которые охватывают ключевые аспекты работы системы: доступность компонентов, использование вычислительных ресурсов, состояние базы данных и эффективность выполнения запросов.

\begin{figure}[ht!] 
	\center
	\includegraphics [scale=0.15] {/Users/vadimvasilev/Study/8 семестр/НИР/latex/SPbPU-student-thesis-template/my_folder/images//7}
	\caption{Список правил формирования оповещений} 
	\label{fig:allalerts}  
\end{figure}

\newpage

С точки зрения логической структуры, все оповещения можно разделить на несколько групп:

\begin{itemize}
	\item Оповещения доступности компонентов — контролируют состояние системы мониторинга и СУБД. К ним относятся \textit{ExporterDown} и \textit{GreenplumDown}, которые фиксируют недоступность соответствующих сервисов.
	
	\item Оповещения использования системных ресурсов — отслеживают нагрузку на инфраструктуру, включая процессор, оперативную память и дисковое пространство. В данную группу входят \textit{HighCPUUsage}, \textit{HighMemoryUsage}, \textit{LowDiskSpace} и \textit{HighDiskIOTime}.
	
	\item Оповещения состояния базы данных — характеризуют активность и поведение СУБД, включая количество подключений и транзакционную активность. К ним относятся \textit{TooManyConnections}, \textit{TooManyIdleConnections} и \textit{NoDBActivity}.
	
	\item Оповещения производительности и качества выполнения запросов — позволяют выявлять потенциальные проблемы в работе запросов и транзакций. В данную группу входят \textit{HighRollbackRate} и \textit{LongRunningQueries}.
\end{itemize}

Каждое правило оповещения содержит три основных параметра: условие срабатывания (expr), время устойчивого нарушения условия (for) и уровень критичности (labels.severity). Использование параметра \textit{for} позволяет исключить кратковременные всплески метрик и формировать уведомления только при устойчивом отклонении значений.

Примером может служить правило контроля загрузки процессора (\firef{fig:alertexample}), при котором формируется предупреждение, если использование CPU превышает 85\% в течение более чем 2 минут. Аналогично, при превышении порога использования оперативной памяти выше 90\% формируется критическое уведомление.

\begin{figure}[ht!] 
	\center
	\includegraphics [scale=0.15] {/Users/vadimvasilev/Study/8 семестр/НИР/latex/SPbPU-student-thesis-template/my_folder/images//8}
	\caption{Примеры правил для оповещений \textit{HighCPUUsage}, \textit{HighMemoryUsage}} 
	\label{fig:alertexample}  
\end{figure}

Alertmanager выполняет маршрутизацию уведомлений на основе заданных правил. В данной работе настроена отправка сообщений на электронную почту, что обеспечивает централизованный контроль состояния системы. Также используется механизм группировки оповещений, который объединяет связанные события и снижает количество отдельных уведомлений.

Таким образом, разработанная система оповещений обеспечивает автоматизированный контроль состояния СУБД Greenplum, позволяя своевременно выявлять критические и потенциально опасные ситуации, связанные с доступностью компонентов, нагрузкой на систему и эффективностью выполнения запросов. Это повышает общую надёжность и устойчивость функционирования системы мониторинга.

\section{Практическое использование разработанной системы} \label{ch2:usage}

Разработанная система мониторинга ориентирована на практическое применение и предполагает минимальное участие пользователя в процессе настройки и развертывания. При проектировании особое внимание уделялось автоматизации конфигурации и упрощению запуска системы.

Процесс использования системы включает несколько последовательных этапов. На первом этапе пользователю необходимо заполнить файл переменных окружения \texttt{.env}, в котором задаются основные параметры конфигурации. К ним относятся данные для подключения к СУБД Greenplum (адрес, порт, учетные данные), параметры сетевых портов сервисов, настройки доступа к Grafana, а также параметры почтового сервера для отправки уведомлений.

Пример структуры файла \texttt{.env} приведён в приложении ~\ref{lst:envexample}:

После заполнения файла переменных окружения выполняется запуск скрипта генерации конфигурационных файлов (~\ref{lst:generateconfigs}):

\begin{verbatim}
	./scripts/generate-configs.sh
\end{verbatim}

Данный скрипт автоматически формирует необходимые конфигурационные файлы для всех компонентов системы, включая Prometheus, Alertmanager и экспортеры, на основе заданных параметров. Это позволяет исключить ручное редактирование конфигураций и снизить вероятность ошибок при настройке.

На завершающем этапе осуществляется запуск системы с использованием Docker Compose:

\begin{verbatim}
	docker compose up -d
\end{verbatim}

Содержание файла можно посмотреть в приложении ~\ref{lst:dockercompose}.

В результате выполнения данной команды происходит автоматическое развертывание всех компонентов системы и их инициализация. После запуска система становится полностью готовой к работе без необходимости дополнительной настройки.

Важной особенностью разработанного решения является высокий уровень автоматизации: пользователю достаточно задать исходные параметры и выполнить две команды, после чего он получает функционирующую систему мониторинга. Такой подход существенно упрощает процесс внедрения и делает систему доступной для использования без глубоких знаний в области настройки отдельных компонентов.

Практическая ценность разработанной системы заключается в том, что она представляет собой не набор разрозненных инструментов, а целостное интегрированное решение, обеспечивающее полный цикл мониторинга — от сбора метрик до генерации и доставки оповещений. Это позволяет сократить время развертывания, снизить сложность эксплуатации и повысить надёжность мониторинга СУБД Greenplum.

Таким образом, предложенный подход обеспечивает удобство использования, воспроизводимость развертывания и готовность системы к применению в реальных условиях.


\section{Выводы по главе} \label{ch2:conclusion}

В рамках данной главы была рассмотрена реализация системы мониторинга СУБД Greenplum, включающей сбор, хранение, визуализацию метрик и механизм оповещений.

Была разработана архитектура системы, основанная на использовании контейнеризации и инструментов Docker Compose, что позволило обеспечить изоляцию компонентов, упростить развертывание и повысить воспроизводимость решения. Все элементы системы интегрированы в единую инфраструктуру и взаимодействуют между собой автоматически.

Реализован механизм централизованного сбора метрик с использованием Prometheus и набора специализированных экспортеров, обеспечивающих получение как системных, так и прикладных показателей работы СУБД. Для анализа состояния системы разработан пользовательский дашборд в Grafana, включающий несколько логических блоков, отражающих различные аспекты функционирования базы данных и инфраструктуры.

Дополнительно реализована система оповещений на основе Alertmanager, включающая набор правил, позволяющих выявлять критические и потенциально опасные ситуации, связанные с доступностью компонентов, нагрузкой на систему и поведением базы данных.

Таким образом, в результате выполнения работы была создана не только архитектура системы мониторинга, но и её полноценная реализация в виде готового программного решения, обеспечивающего комплексный мониторинг СУБД Greenplum и пригодного для практического использования.

